% ============================================================================
% 00_resumo_abstract.tex
% Substitui o RESUMO (pp. 7) e o ABSTRACT (pp. 8) do TCC original.
%
% Como usar:
%   1. Localize no .tex original os ambientes \begin{resumo} ... \end{resumo}
%      e \begin{abstract} ... \end{abstract} (ou equivalentes do abntex2).
%   2. Substitua o conteudo interno por este arquivo.
% ============================================================================

\begin{resumo}
O crescimento da temperatura global é um fato alarmante. Segundo a \emph{World Meteorological Organization}~\cite{wmo2021}, o ano de 2020 atingiu uma média histórica de \SI{1,2}{\degree C} acima das temperaturas aferidas na era pré-industrial. Aragão \emph{et al.}~\cite{aragao2018amazon} documentam que, desde o início do século~XXI, o fogo associado à seca passou a \emph{neutralizar} os ganhos obtidos com a redução do desmatamento na Amazônia; Silva-Junior \emph{et al.}~\cite{silvajunior2025amazon} mostram que 2024 marcou a pior degradação por fogo em mais de duas décadas --- \SI{3,3}{\mega\hectare} queimados (\(+400\%\) em relação aos 2 anos anteriores), com emissão estimada de \SI{791}{\mega\tonne} de \ce{CO2}. A preservação das áreas verdes é crucial para mitigar este cenário, e ferramentas de previsão automatizadas, baseadas em aprendizado de máquina, representam uma abordagem promissora.

Este projeto propõe um sistema de classificação supervisionada de risco de incêndio na Amazônia Legal em três níveis (Baixo, Moderado, Muito Alto), treinado sobre \num{924306} registros do Programa Queimadas (2014--2023). O \emph{pipeline} integra \textbf{24~\emph{features} físico-climáticas Tier~1} (SPI, KBDI \emph{proxy}, VPD \emph{proxy}, anomalias de precipitação, médias móveis estendidas, histórico estendido de fogo), substituindo metodologicamente a \emph{umidade relativa} --- cuja integração via NASA POWER demandaria 16 dias contínuos de enriquecimento --- por \emph{proxies} amparados pela literatura recente~\cite{seager2015climatology,forests2024cerradoamazon,npjhazards2025fire}. Foram comparadas oito abordagens de aprendizado supervisionado (Regressão Logística, SGD, Random Forest com SMOTE, RF balanceado otimizado por \textbf{Optuna}~\cite{akiba2019optuna}, XGBoost otimizado~\cite{chen2016xgboost}, LightGBM otimizado~\cite{ke2017lightgbm}, CatBoost~\cite{prokhorenkova2018catboost}, \emph{ensembles} \emph{Voting} e \emph{Stacking}~\cite{wolpert1992stacked}), com calibração isotônica~\cite{zadrozny2002transforming} e ajuste multi-classe de limiares de decisão.

O melhor modelo --- \textbf{\emph{Ensemble Stacking}} com meta-classificador \emph{Gradient Boosting}~\cite{friedman2001greedy} (\texttt{ensemble\_stacking\_gbm}) --- atingiu \textbf{\SI{84,61}{\percent}} de acurácia, \textbf{0,7995} de \(\text{F}_1\)-macro e \textbf{0,6375} de \(\text{F}_1\)-Moderado (classe historicamente mais difícil) em \emph{hold-out} estratificado de \num{184862} amostras --- um ganho de \(+4,12\) pontos percentuais (pp) em acurácia e \(+8,77\,\text{pp}\) em \(\text{F}_1\)-Moderado sobre o \emph{ensemble} \emph{baseline}. A interpretabilidade é abordada em dupla técnica: \textbf{SHAP por classe}~\cite{lundberg2017unified,lundberg2020local} via \emph{Random Forest} compacto \emph{proxy} e \textbf{\emph{Permutation Importance}} \emph{model-agnostic}~\cite{breiman2001random,fisher2019all} sobre o \emph{Stacking} GBM final --- ambas confirmam \texttt{KBDI\_proxy}, \texttt{VPD\_proxy}, \texttt{Indice\_Seca}, \texttt{DiaSemChuva\_ma90} e \texttt{Precipitacao\_acum\_180d} como variáveis dominantes. Uma \textbf{validação temporal estrita} (\emph{rolling-origin} por ano, 2021--2023)~\cite{bergmeir2012cv} quantifica o viés do \emph{split} aleatório frente a uma avaliação causal (queda de \SI{14,49}{pp} em acurácia, \SI{32,35}{pp} em \(\text{F}_1\)-Moderado). O modelo final é exposto via aplicação web \emph{Flask} + \emph{Folium}~\cite{folium,leafletjs} com explicação local e integração de clima em tempo real (NASA POWER~\cite{nasapower}) e detecção ativa de focos (NASA FIRMS~\cite{nasafirms}).

\textbf{Palavras-chave}: Aprendizado de Máquina. Previsão. Queimadas. Amazônia Legal. Ensemble Stacking. XGBoost. SHAP. Interpretabilidade.
\end{resumo}

% ----------------------------------------------------------------------------

\begin{abstract}
The alarming rise in global temperatures is an undeniable reality. According to the World Meteorological Organization~\cite{wmo2021}, the year 2020 witnessed a historical average temperature increase of \SI{1.2}{\degree C} above pre-industrial levels. Aragão \emph{et al.}~\cite{aragao2018amazon} document that, since the beginning of the 21st century, drought-related fires have started to \emph{offset} the gains obtained from reductions in Amazon deforestation; Silva-Junior \emph{et al.}~\cite{silvajunior2025amazon} report that 2024 marked the worst fire-driven degradation event in over two decades --- \SI{3.3}{\mega\hectare} burned (\(+400\%\) over the previous two years), with estimated emissions of \SI{791}{\mega\tonne} of \ce{CO2}. Preserving forested areas is therefore crucial, and automated prediction tools based on machine learning represent a promising approach.

This project proposes a supervised classification system for fire risk in the Brazilian Legal Amazon with three severity levels (Low, Moderate, Very High), trained on \num{924306} records from the INPE Fire Programme (2014--2023). The pipeline integrates \textbf{24~Tier~1 physico-climatic features} (SPI, KBDI proxy, VPD proxy, precipitation anomalies, extended moving averages, extended fire history), methodologically replacing the \emph{relative humidity} variable --- whose integration via NASA POWER would require 16 continuous days of enrichment --- with proxies supported by recent literature~\cite{seager2015climatology,forests2024cerradoamazon,npjhazards2025fire}. Eight supervised learning approaches were compared (Logistic Regression, SGD, Random Forest with SMOTE, balanced RF optimised with \textbf{Optuna}~\cite{akiba2019optuna}, optimised XGBoost~\cite{chen2016xgboost}, optimised LightGBM~\cite{ke2017lightgbm}, CatBoost~\cite{prokhorenkova2018catboost}, \emph{Voting} and \emph{Stacking}~\cite{wolpert1992stacked} ensembles), with isotonic probability calibration~\cite{zadrozny2002transforming} and multi-class decision-threshold tuning.

The best model --- a \textbf{\emph{Stacking Ensemble}} with a \emph{Gradient Boosting} meta-classifier~\cite{friedman2001greedy} --- reached \textbf{\SI{84.61}{\percent}} accuracy, \textbf{0.7995} macro-\(\text{F}_1\) and \textbf{0.6375} \(\text{F}_1\) on the Moderate class (historically the hardest) in a stratified \emph{hold-out} of \num{184862} samples, a \(+4.12\) percentage-point gain in accuracy and \(+8.77\,\text{pp}\) in Moderate-\(\text{F}_1\) over the \emph{ensemble} baseline. Interpretability is addressed by two complementary techniques: class-wise \textbf{SHAP}~\cite{lundberg2017unified,lundberg2020local} via a compact Random Forest proxy and model-agnostic \textbf{Permutation Importance}~\cite{breiman2001random,fisher2019all} on the final Stacking GBM --- both identify \texttt{KBDI\_proxy}, \texttt{VPD\_proxy}, \texttt{Indice\_Seca}, \texttt{DiaSemChuva\_ma90} and \texttt{Precipitacao\_acum\_180d} as dominant variables. A strict \textbf{temporal validation} (yearly rolling-origin, 2021--2023)~\cite{bergmeir2012cv} quantifies the bias of the random split against a causal evaluation (drop of \SI{14.49}{pp} in accuracy and \SI{32.35}{pp} in Moderate-\(\text{F}_1\)). The final model is exposed through a \emph{Flask} + \emph{Folium}~\cite{folium,leafletjs} web application with local explanation and real-time climate (NASA POWER~\cite{nasapower}) and active fire detection (NASA FIRMS~\cite{nasafirms}) integration.

\textbf{Keywords}: Machine Learning. Prediction. Wildfires. Legal Amazon. Stacking Ensemble. XGBoost. SHAP. Interpretability.
\end{abstract}
